%! AP Statistics — Unit 7, Lesson 7.2 — Homework
\documentclass[10pt]{article}
\usepackage{apstats-article}
\usepackage{apstats-boxes}
\newcommand{\TallMath}[1]{$\displaystyle #1\rule[-1.4em]{0pt}{3.2em}$}

\begin{document}

\pageheader{Unit 7, Lesson 7.2}{Homework}
\namedateperiod

\vspace{0.1in}
\textbf{Directions:} Use the full four-step procedure for each confidence interval problem.

\begin{enumerate}

  \item \textbf{Screen Time.} A pediatrician randomly samples 20 elementary school
        students and records their average daily screen time. Summary statistics:
        $\bar x = 3.4$ hours, $s = 0.9$ hours.

        Construct a 95\% confidence interval for the true mean daily screen time
        for elementary school students. Verify all conditions.
        \vspace{0.1in}

        \textbf{Step 1:} \writeline
        \textbf{Step 2 (Conditions):} \writelines{3}
        \textbf{Step 3:} $df =$ \blank{2cm}, $t^* =$ \blank{2.5cm}, $SE =$ \blank{2.5cm}

        Interval: \blank{8cm} $= ($ \blank{2cm}$,$ \blank{2cm}$)$ hours
        \vspace{0.05in}

        \textbf{Step 4:} \writelines{2}

  \item \textbf{Fuel Efficiency.} A car manufacturer randomly selects 40 vehicles of a
        new model and measures fuel efficiency: $\bar x = 31.7$ mpg, $s = 4.2$ mpg.

        Construct a 99\% confidence interval for the true mean fuel efficiency. Verify
        conditions and interpret in context.
        \vspace{0.1in}

        \textbf{Steps 1--4:} \writelines{8}

        Interval: $= ($ \blank{2cm}$,$ \blank{2cm}$)$ mpg

  \item \textbf{Matched Pairs: Reaction Time.} A researcher tests whether caffeine
        affects reaction time. Twelve participants are tested before and after consuming
        caffeine. The differences in reaction time (before $-$ after, in milliseconds) are:

        \medskip
        \begin{center}
        \begin{tabular}{cccccccccccc}
            14 & $-3$ & 22 & 8 & 11 & $-5$ & 18 & 6 & 13 & 20 & $-2$ & 15
        \end{tabular}
        \end{center}
        \medskip

        Compute $\bar d$ and $s_d$, then construct a 90\% confidence interval for
        the mean change in reaction time due to caffeine.

        $\bar d = $ \blank{3cm} \quad $s_d \approx $ \blank{3cm}

        \textbf{Steps 1--4:} \writelines{6}

        Interval: $= ($ \blank{2cm}$,$ \blank{2cm}$)$ ms

\end{enumerate}

\begin{extensionbox}
\textbf{Extension --- optional.} For Problem 1, suppose $\sigma$ were known to be
exactly 0.9 hours. Compute the $z$-interval for the same data and confidence level.
Compare the width of the $z$-interval to the $t$-interval. Explain the difference.

\writelines{3}
\end{extensionbox}

\end{document}
